\section{Conclusion}

The narrow claim of this paper is that agent coordination topology constrains what workers and
coordinators can observe, and that those visibility constraints are strong enough to predict several
practically relevant tradeoffs: when review helps, when decomposition hurts, when parallelism pays off,
and when tool fragmentation becomes a tax. The formal apparatus is useful to the extent that it sharpens
those architecture choices.

The practical implication is that topology should be treated as a first-class design variable in agent
engineering. Single agents remain strong on strictly sequential tasks. Reviewer or hub structures help
when multiple workers can independently introduce costly errors. Specialist swarms pay off when subtasks
are close to epistemically independent and coordinator overhead is small. Tool-heavy workloads punish
gratuitous fragmentation unless routing is itself part of the design.

This paper is therefore not a claim that one topology dominates all others, nor a claim that the current
theory replaces benchmark evaluation. It is a step toward a more explicit architecture calculus for AI
agent systems: one in which coordination choices are justified by visibility, decomposition structure, and
cost rather than by orchestration folklore alone.
